

%%% Choose a language %%%

\newif\ifEN
\ENtrue   % uncomment this for english
%\ENfalse   % uncomment this for czech

%%% Configuration of the title page %%%

\def\ThesisTitleStyle{mff} % MFF style
%\def\ThesisTitleStyle{cuni} % uncomment for old-style with cuni.cz logo
%\def\ThesisTitleStyle{natur} % uncomment for nature faculty logo

\def\UKFaculty{Faculty of Mathematics and Physics}
%\def\UKFaculty{Faculty of Science}

\def\UKName{Charles University in Prague} % this is not used in the "mff" style

% Thesis type names, as used in several places in the title
%\def\ThesisTypeTitle{\ifEN BACHELOR THESIS \else BAKALÁŘSKÁ PRÁCE \fi}
\def\ThesisTypeTitle{\ifEN MASTER THESIS \else DIPLOMOVÁ PRÁCE \fi}
%\def\ThesisTypeTitle{\ifEN RIGOROUS THESIS \else RIGORÓZNÍ PRÁCE \fi}
%\def\ThesisTypeTitle{\ifEN DOCTORAL THESIS \else DISERTAČNÍ PRÁCE \fi}
%\def\ThesisGenitive{\ifEN bachelor \else bakalářské \fi}
\def\ThesisGenitive{\ifEN master \else diplomové \fi}
%\def\ThesisGenitive{\ifEN rigorous \else rigorózní \fi}
%\def\ThesisGenitive{\ifEN doctoral \else disertační \fi}
%\def\ThesisAccusative{\ifEN bachelor \else bakalářskou \fi}
\def\ThesisAccusative{\ifEN master \else diplomovou \fi}
%\def\ThesisAccusative{\ifEN rigorous \else rigorózní \fi}
%\def\ThesisAccusative{\ifEN doctoral \else disertační \fi}



%%% Fill in your details %%%

% (Note: \xxx is a "ToDo label" which makes the unfilled visible. Remove it.)
\def\ThesisTitle{Effect of regularization on classification of magnetic topological phases}
\def\ThesisAuthor{Andrej Rendek}
\def\YearSubmitted{2025}

% department assigned to the thesis
\def\Department{Department of Condensed Matter Physics}
% Is it a department (katedra), or an institute (ústav)?
\def\DeptType{Department}

\def\Supervisor{RNDr. Pavel Baláž, Ph.D.}
\def\SupervisorsDepartment{Department of Condensed Matter Physics}

% Study programme and specialization
\def\StudyProgramme{Mathematical and Computational Modelling in Physics}
\def\StudyBranch{FMPMP}

\def\Dedication{%
Dedication. I would like to thank my advisor RNDr. Pavel Baláž, Ph.D. for his guidance, for introducing me to the field of machine learning, and for his patience throughout this journey.

I also want to thank my parents for their support and Natalia for all the love and warmth you always bring me, both in joyful and more gloomy times.
}

\def\AbstractEN{%
In this thesis we study the implementation of neural networks for the problem of classification of simulated magnetic configurations into 3  topological phases: ferromagnet, magnetic skyrmion and magnetic spiral. The main challenge lies in the fact that the phase transitions are continuous, meaning two phases can occur in the same configuration. Even though we don't have true labels for such configurations, we aim to improve the models' performance exactly on such cases. To achieve that, we explore different model architectures, including dense and convolutional networks, choice of hyperparameters including activation functions, and the effect of various regularization techniques such as Dropout Monte Carlo or data augmentations. We found out that simple and quite shallow, convolution networks without attention mechanism work the best, and can reliably classify the single phase configurations, while performing better on the phase transitions compared to the deeper networks. The CutMix data augmentation improves performance on the skyrmion-spiral transition by about $60\%$ according to by us defined metric, compared to the results without augmentations. At the same time, CutMixes harm the perfomance on ferromagnetic-skyrmion boundary. MixUp provides small, but complementary gains. By applying these augmentations, the model learns to predict the probability of a skyrmion phase during phase transition into spiral that alighns with the change in topological charge. On the other hand, Monte Carlo dropout shows no benefit.
}


\def\AbstractCS{%
V této práci se zabýváme implementací neuronových sítí pro klasifikaci simulovaných magnetických konfigurací do tří topologických fází: feromagnet, magnetický skyrmion a magnetická spirála. Hlavní výzva spočívá ve spojitých fázových přechodech, což znamená, že se ve stejné konfiguraci mohou vyskytovat dvě fáze. Přestože pro takové konfigurace nemáme k dispozici lable, snažíme se zlepšit predikce modelů právě v takových případech. Abychom toho dosáhli, zkoumáme různé architektury modelů, včetně plně propojených a konvolučních sítí, volbu hyperparametrů, včetně aktivačních funkcí, a vliv různých regularizačních technik, jako je Dropout Monte Carlo nebo augmentace dat. Zjistili jsme, že nejlépe fungují jednoduché, mělké konvoluční sítě. Dokážou spolehlivě klasifikovat jednofázové konfigurace a zároveň si vedou lépe než hlubší sítě při fázových přechodech. Augmentace dat CutMix zlepšuje výkon na skyrmionově spirálním přechodu podle námi definované metriky přibližně o $60\%$ ve srovnání s výsledky bez augmentací. CutMixy zároveň poškozují predikce na hranici feromagnetu a skyrmionu. MixUp poskytuje malé, ale doplňkové zlepšení. Použitím těchto augmentací se model naučí předpovídat pravděpodobnost skyrmionové fáze při fázovém přechodu do spirály způsobem, který je v souladu se změnou topologického náboje v těchto konfiguracích. Na druhou stranu regularizace dropout Monte Carlo nevykazuje žádný přínos.
}

% 3 to 5 keywords (recommended), each enclosed in curly braces.
% Keywords are useful for indexing and searching for the theses by topic.
\def\Keywords{%
{machine learning} {magnetic skyrmion} {neural networks}
}

% If your abstracts are long and do not fit in the infopage, you can make the
% fonts a bit smaller by this setting. (Also, you should try to compress your abstract more.)
% Alternatively, consider increasing the size of the page by uncommenting the
% geometry modification in thesis.tex.
\def\InfoPageFont{}
%\def\InfoPageFont{\small}  %uncomment to decrease font size

\ifEN\relax\else
% If you are writing a czech thesis, you additionally need to fill in the
% english translation of the metadata here!
\def\ThesisTitleEN{\xxx{Thesis title in English}}
\def\DepartmentEN{\xxx{Name of the department in English}}
\def\DeptTypeEN{\xxx{Department}}
\def\SupervisorsDepartmentEN{\xxx{Superdepartment}}
\def\StudyProgrammeEN{\xxx{study programme}}
\def\StudyBranchEN{\xxx{study branch}}
\def\KeywordsEN{%
\xxx{{key} {words}}
}
\fi
